\textbf{Lo que muestran julio y agosto de 2026.} El sistema se calibró con la G-21 y se probó en la 60-CH. En la G-21 funciona: reproduce el estado observado con kappa ponderado de 0,54, contra 0,27 del cuadro del protocolo de Lo Barnechea, y clasifica como cierre por peligro 10 de los 11 bloques del corte del 17 al 20 de julio sin ver el estado observado. Pero lo hace por el peligro de avalancha del tramo alto, mientras que el corte se decretó por rodados, nieve y un aluvión del río Ñilhue en el tramo medio, que el índice de remoción en masa no detectó. La regla de Walker pura se ajusta mejor a la práctica (0,63), porque el piso del protocolo de Lo Barnechea manda a Alerta 2 la ``presencia de lluvia'' con isoterma alta y exige cierre total con poco más de 1 cm de nieve cuando la isoterma baja de 1.500 m, y la práctica no hace ninguna de las dos cosas. El peligro, en cambio, tiene que ir aparte: metido como una variable más, el sistema cierra la 60-CH en 1 de sus 176 bloques cerrados.

\textbf{El cierre largo de la 60-CH.} De los 36 días cerrados entre el 14 de julio y el 18 de agosto, el sistema respalda cierre total en 20, casi todos por peligro de avalancha. En los otros 16 justifica cierre parcial (3), restricción (9) o nada (4). Contando todo el periodo, entre el 41\,\% y el 58\,\% de las 1.019 horas cerradas no tienen respaldo meteorológico para cierre total, según la capacidad de despeje. Esto no prueba que la decisión haya sido errada. Primero, el resultado depende de la exposición del tramo alto a avalanchas, que hoy es un supuesto. Segundo, el sistema no ve la visibilidad: del 5 al 8 de agosto hubo ráfagas de hasta 31 m/s. Lo que sí muestra es que había opciones intermedias, y qué dato falta para elegir entre ellas: un catastro de trayectorias de avalancha.
